\documentclass[10pt,letterpaper,sans]{moderncv}
\moderncvstyle{banking}
\moderncvcolor{blue}

\usepackage[scale=0.84]{geometry}
\nopagenumbers{}
\setlength{\hintscolumnwidth}{2.2cm}

\name{Srikrishnaa}{Vadivel}
\title{Computational Physicist | First Principles | HPC | Scientific Software}
\email{srikrishnaacareers@gmail.com}
\phone[mobile]{516-853-5663}
\social[github]{krishnaa423}
\extrainfo{\href{https://leetcode.com/u/krishnaa42342}{leetcode.com/u/krishnaa42342} \textbar{} \href{https://hub.docker.com/u/krishnaa42342}{hub.docker.com/u/krishnaa42342}}

\begin{document}
\makecvtitle

\section{Profile}
\cvitem{}{Computational physics PhD researcher focused on first-principles calculations of self-trapped excitons, scientific software, and high-performance simulation. Experience includes leadership-class HPC, PETSc/SLEPc, MPI/OpenMP, Python/C++, and ML-assisted materials workflows.}

\section{Research}
\cventry{2021--present}{PhD Researcher}{Yale University}{}{}{\begin{itemize}%
\item Perform first-principles calculations of self-trapped excitons using HPC resources at Oak Ridge, Aurora, LBNL, and LLNL.
\item Build reproducible Linux, Bash, Python, C/C++, CMake, PETSc/SLEPc, MPI/OpenMP, Docker, and Kubernetes workflows for large simulations.
\item Explore machine-learning surrogates and equivariant molecular models for materials and molecular systems.
\end{itemize}}

\section{Selected Technical Projects}
\cvitem{DOLFINx PDE Gallery}{Weak forms for Poisson, heat, and elasticity, plus solver and visualization scaffold.}
\cvitem{First-Principles LiF}{Plane-wave pseudopotential Hamiltonian and band plotting.}
\cvitem{Physics Action Notes}{Path integrals, imaginary time, EM, GR, QED/QCD, Standard Model, and strings from action principles.}
\cvitem{Meep Ring Resonator}{Photonics FDTD geometry and source setup.}

\section{Industry}
\cventry{2019--2021}{Software and Embedded Engineer}{Probe Technologies / Weatherford}{}{}{Built CUDA/OpenGL waveform processing and visualization software for acoustic logging tools; also delivered embedded sensor logging firmware and FPGA logic.}

\section{Education}
\cventry{2021--present}{PhD, Electrical Engineering}{Yale University}{}{}{}
\cventry{2023}{MPhil, Electrical Engineering}{Yale University}{}{}{}
\cventry{2023}{MS, Electrical Engineering}{Yale University}{}{}{}
\cventry{2017--2018}{MEng, Electrical and Computer Engineering}{Cornell University}{}{}{}
\cventry{2013--2017}{BA, Electrical and Computer Engineering}{Cornell University}{}{}{}

\section{Skills}
\cvitem{Methods}{DFT workflows, PDEs, numerical linear algebra}
\cvitem{Scientific}{PETSc, SLEPc, DOLFINx concepts, Meep concepts}
\cvitem{Systems}{MPI, OpenMP, Linux, CMake}
\cvitem{Programming}{Python, C++, CUDA, PyTorch}

\end{document}
